\newcommand{\DoNotLoadEpstopdf}{}
\documentclass[11pt,reqno]{amsart}
\usepackage[T1]{fontenc}
\usepackage{lmodern,amsmath,amssymb,amsthm,mathtools,mathrsfs,microtype}
\usepackage[margin=1.1in]{geometry}
\usepackage[hidelinks]{hyperref}
\hypersetup{pdftitle={Collision records, arithmetic, and raw local inversion in a triangular Lorentz gas},pdfauthor={Qian Qi}}
\setlength{\emergencystretch}{2em}
\allowdisplaybreaks[2]
\numberwithin{equation}{section}
\newtheorem{theorem}{Theorem}[section]
\newtheorem{lemma}[theorem]{Lemma}
\newtheorem{proposition}[theorem]{Proposition}
\newtheorem{corollary}[theorem]{Corollary}
\theoremstyle{definition}
\newtheorem{definition}[theorem]{Definition}
\newtheorem{inputstatement}[theorem]{Input statement}
\theoremstyle{remark}\newtheorem{remark}[theorem]{Remark}
\newcommand{\R}{\mathbb R}\newcommand{\Z}{\mathbb Z}
\newcommand{\one}{\mathbf1}\newcommand{\dd}{\,\mathrm d}
\newcommand{\esssup}{\operatorname*{ess\,sup}}
\newcommand{\diver}{\operatorname{div}}
\newcommand{\TV}{\operatorname{TV}}
\title[Collision records and raw local inversion]{Collision records, arithmetic, and raw local inversion in a triangular Lorentz gas}
\author{Qian Qi}
\date{October 11, 2026. A2-DYN, revision 78}
\subjclass[2020]{37D50, 37A50, 60F05, 28A75}
\keywords{Sinai billiards, exact collision records, local limits, endpoint momenta, complete-source transport, occupation boundary currents}
\begin{document}
\raggedbottom
\begin{abstract}
We study the pointwise local limit problem for the complete displacement,
collision-count, occupation and roof record of actual returns in a
triangular Lorentz gas. The physical endpoint-momentum field gives
collision-uniform control of every contact, with an exponential loss
of influence toward the middle of the word. This holds on all regular
surviving trajectories, without a lower incidence cutoff. It yields a
uniform square-root modulus for the original full guard. By transporting
its error to the target fiber and averaging before the collision limit,
we estimate the complete guard-distortion contribution at the required
normalization using only fixed-width concentration. The remaining
geometric contribution retains true exits and physical coarea-volume
loss. We also prove an exact discrete-divergence identity for internal
section-decision currents, with a peak-aware Fourier estimate that keeps
nontrivial arithmetic phases. These results advance the original
pointwise theorem without replacing its exact labels or its source;
the complete true-exit and geometric-distortion estimate is the
quantitative step still to be established.
\end{abstract}
\maketitle

\section{The unchanged physical problem}
Put $R_-=9/20$ and $R_+=47/100$, and let $R\in I=[R_-,R_+]$. Circular scatterers have centers in
$\Lambda=\Z(1,0)+\Z(1/2,\sqrt3/2)$. Write $T_R$ for the regular
collision map and $(\alpha,p)\in(\R/2\pi\Z)\times(-1,1)$ for its
coordinates; $p$ is the outgoing tangential velocity and
$\dd\nu=\dd\alpha\,\dd p/(4\pi)$. Physical flights satisfy
$\tau_-\le\tau_R\le\tau_+$, where $\tau_-=3/50$ and $\tau_+$
is the fixed finite-horizon bound.

The section is specified here to fix every event in the argument.
Put $h_0=1/200$. Let $Y_R$ be the union of the four open coordinate
squares of half-side $h_0$ centered at
\[
 (0,0),\quad(\pi/6,-1/2),\quad
 (-\theta_R,\sin\theta_R),\quad
 (\pi/3-\theta_R,\sin\theta_R).
\]
Here $\theta_R=2\arctan q_R$, and $q_R$ is the unique root in
$(49/100,53/100)$ of
\[
 8Rq(1+q^2)-4q(1-q^2)+\sqrt3(1-6q^2+q^4)=0.
\]
Set
\[
 U=\{|\alpha-\pi/6|<3/50,\ |p|<3/20\},\qquad
 Y_R^*=Y_R\cup U,\qquad c_*=\frac{91}{10000\pi}.
\]
The rectangles are disjoint and $\nu(Y_R^*)=c_*$. We always use
$\nu_R^*=c_*^{-1}\nu|_{Y_R^*}$, not a renormalization of a smaller
source.

At $m$ flights, let $L=L_{m,R}$ be their total length,
$K^{\rm c}_{m,R}$ their lattice displacement, and
$A_{m,R}=\sum_{j=0}^{m-1}\one_{Y_R^*}\circ T_R^j$. The event is
\begin{equation}\label{eq:v76-event}
 \Pi_{n,k,m,R}=Y_R^*\cap\{K^{\rm c}_{m,R}=k,\ A_{m,R}=n\}
                         \cap T_R^{-m}Y_R^*.
\end{equation}
The initial visit is counted; the terminal visit is tested separately.
On this event the actual $n$th return has collision count $m$ and
roof $L$. Denote its density by $p_{n,R}(k,m,t)$.

Write $c_j=\sqrt{1-p_j^2}$ and let $\Delta_j$ be the least clearance
of flight $j$ from any third disk, excluding its two incident disks.
Keep the original smooth function $0\le\Theta\le1$, zero on
$(-\infty,1]$ and one on $[2,\infty)$. With $q_0=47/53$,
$d_j=\min(j,m-j)$ and $d'_j=\min(j,m-1-j)$, set
\begin{align}
 H^\varepsilon&=
 \prod_{j=0}^{m}\Theta\!\left(\frac{c_j}{\varepsilon q_0^{d_j/4}}\right)
 \prod_{j=0}^{m-1}\Theta\!\left(\frac{\Delta_j}{\varepsilon q_0^{d'_j/4}}\right),
 \qquad W^\varepsilon=1-H^\varepsilon,\label{eq:v76-guard}\\
 b^\varepsilon(t)\dd t
   &=L_\#(W^\varepsilon\one_{\Pi_{n,k,m,R}}\nu_R^*).
 \label{eq:v76-source}
\end{align}
In particular first incidence and first clearance outside any selected
critical neighborhoods are still present. No section-distance factor
has been inserted into this guard.

\subsection{The arithmetic reference and the target}
For $y=(k_1,k_2,t,n)$ put
\[
 \mu_R=(0,0,\bar\tau_R,c_*),\quad
 Z=(y-m\mu_R)/\sqrt m,\quad
 \mathcal C_{m,R,n,k}(M)=\{t:|Z|\le M\}.
\]
We retain the finite spectral-peak reference, including peaks close to
but not on the unit circle. To specify its meaning, the collision
record observable is $(\kappa_1,\kappa_2,\tau,\one_{Y_R^*})$.
Let $\lambda_j(R,\vartheta)$ and $\mathsf P_j(R,\vartheta)$ be its
continued simple twisted eigenvalue and projection in each of the
fixed peak charts, and let $\chi_j$ be their frequency partition.
At the real maximum $a_j$ of $\log|\lambda_j|$ put
\[
 z_j=\lambda_j(R,a_j),\quad
 \mu_j=i^{-1}\nabla\log\lambda_j(R,a_j),\quad
 H_j=-\nabla^2\log\lambda_j(R,a_j),
\]
with $\operatorname{Re}H_j\ge aI_4$ for a common $a>0$. If
$\eta_R=\one_{Y_R^*}$ and $\ell$ denotes mass evaluation, put
$B_j=\chi_j(R,a_j)\ell(M_{\eta_R}\mathsf P_j(R,a_j)(\eta_R\nu))$.
Define
\begin{align}
 G_H(z)&=(2\pi)^{-4}\int_{\R^4}e^{-iv\cdot z-v^{\mathsf T}Hv/2}\dd v,
 \notag\\
 \mathcal L_{m,R}(y)&=\operatorname{Re}\left[
 c_*^{-1}\sum_j e^{-ia_j\cdot y}z_j^m B_j
      G_{H_j}\left((y-m\mu_j)/\sqrt m\right)\right].
 \label{eq:v76-reference}
\end{align}
Only the charts active at $R$ are summed. This is the same transition
reference as in the retained spectral theorem. The existence and
uniform estimates for these scalar data belong to that analytic
input, not to the transport argument below. No frequency branch or
arithmetic zero class is deleted.

The governing assertion remains
\begin{equation}\label{eq:v76-target}
 \sup_{R,\,1\le n\le m,\,k\in\Z^2}
 \esssup_{t\in\mathcal C_{m,R,n,k}(M)}
 |m^2p_{n,R}(k,m,t)-\mathcal L_{m,R}(k_1,k_2,t,n)|\longrightarrow0.
\end{equation}
Let $\mathcal E_M$ denote this error's collision limsup. An empty
central target contributes zero. The order of the uniformity is
$\sup_R\esssup_t$, so exceptional density sets may depend on $R$.

\subsection{The analytic inputs in the form used here}
\begin{inputstatement}\label{input:v76-analytic}
The unchanged source \eqref{eq:v76-source} has a nonnegative,
compactly supported $BV$ density at every fixed count. For every
fixed $\varepsilon,\delta>0$, with
$A_\delta f(t)=(2\delta)^{-1}\int_{t-\delta}^{t+\delta}f(s)\dd s$,
\begin{equation}\label{eq:v76-average}
 \limsup_m\sup_{R,n,k}m^2\|A_\delta b^\varepsilon\|_\infty
                                      \le C\varepsilon^{1/16}.
\end{equation}
The constant is independent of the previously fixed width. For large
fixed $B$ and $\varepsilon(B)=A_0B^{-1/12}$,
\begin{equation}\label{eq:v76-remainder}
 m^2p-\mathcal L=m^2b^{\varepsilon(B)}+e_B,\qquad
 \limsup_m\sup_{R,n,k}\esssup_{\mathcal C(M)}|e_B|
                                     \le C_MB^{-1/192}.
\end{equation}
All auxiliary spectral bands in these statements are fixed before
the collision limit and enlarged only afterwards.
\end{inputstatement}
These are the complete-source tame theorem, the fixed-width
complete-source bound, and the positive scalar remainder theorem in
the technical supplement. Their proofs and historical hypotheses are
retained there. The fixed-count source uses threshold integration
\cite{CMInt}, definable monotonicity \cite{Coste}, and the
distributional characterization of $BV$ \cite{AFP}; it does not
assume definability of the original smooth guard. They are restated as inputs to distinguish them from
the new geometric and transport proofs. No long-time bound on a
fixed-count variation constant is part of this statement.


\subsection{The physical construction and the new estimate}
The endpoint-momentum field and complete angular recoupling are
retained below with their proofs. The new argument then estimates
one full dynamical contribution, rather than replacing it by a
per-word bound. All-contact optical response gives a uniform modulus
for the unchanged guard, even with arbitrarily small internal
incidences. Its transported error is measured against a translated
unweighted target; averaging that target permits the fixed-width
analytic input to be used without a pointwise density assumption.

For the remaining physical geometry, internal section interfaces
admit an exact signed difference in the occupation index. That
identity is stated before taking absolute values, and its Fourier
consequence is formulated at every peak rather than only at the
origin. True singular and outer section boundaries remain explicit.

\input{core/174_canonical_endpoint_momentum_flow}
\input{core/175_complete_angular_recoupling}
\input{core/176_angular_current_and_ordered_endpoint}

\input{core/177_all_contact_motion_and_guard}
\input{core/178_averaged_guard_distortion}
\input{core/179_exact_occupation_boundary_current}

\section{The complete geometric estimate and the original theorem}
The pointwise assertion remains \eqref{eq:v76-target}, with the full
arithmetic transition, the same normalization and all zero classes.
Theorem~\ref{thm:v78-separation} proves the complete guard-distortion
bound \eqref{eq:v78-guard-closure}; the constant is uniform in collision
count and no selected physical word has replaced the source.
Proposition~\ref{prop:v78-exit-localization} locates every true forward
exit, but does not bound concentration on the starting roof.
Theorem~\ref{thm:v78-decision-current} computes the signed internal
section contribution without suppressing its exact-label coefficients.

To conclude the organizing theorem, one must still estimate the
complete geometric term in \eqref{eq:v78-ordered-geometry}, or the
complete angular defect in \eqref{eq:v77-ordered-defect}. The new
guard estimate does not control the optical Jacobian loss. An
integrated boundary-layer estimate is not a pointwise source estimate,
and a spectral estimate for volume multipliers is not automatically
an estimate for surface-current insertions. In particular
\eqref{eq:v78-peak-difference} retains nonzero arithmetic leading
terms. The proof ledger records these distinctions explicitly.

The organizing pointwise theorem is therefore not marked proved in
this revision. Its path and likelihood consequences retain the same
scalar-first logical order. The complete inherited mathematical
source and the preceding journal article are retained in the technical
supplement and provenance, rather than deleted or recast as a different
problem.

\input{journal-references}
\end{document}
